\section{The exact polar correction and a pole-free raw diagonal obstruction}
\label{sec:euler-polar-identification}

The multiplier continuation in MA-MO-003 can be compared with the
original summation operator, including its polar term. In the rational
sector that comparison also produces actual smooth inputs on which the
uncontinued first diagonal is unbounded. All assertions below concern
the specified operators; they do not exclude cancellation after
restriction to a smaller common-sum space.

Fix \(0<a<1/2\), \(\sigma=a+1/2\), and a finite set \(S\) of rational
primes. Keep \(\mathcal H_a,V_a\) from (MA-MO.20), and set
\[
 \begin{aligned}
 \mathcal E&=\{U\in C^\infty(\mathbb R):e^{bx}U(x)\in
                    \mathscr S(\mathbb R)\ \hbox{for every }b\in\mathbb R\},\\
 \mathcal E_0&=\{F\in\mathcal E:\mathcal MF(1/2)=0\},\\
 \mathcal MU(z)&=\int_{\mathbb R}e^{zx}U(x)\,dx,\\
 E_S(s)&=\prod_{p\in S}(1-p^{-s}),\qquad
 \kappa_S=E_S(1)>0,\qquad p_0(x)=e^{-x/2}.
 \end{aligned}
 \tag{MA-MO.43}
\]
Here \(E_S(s)\) is the same finite Euler multiplier as in MA-MO.21,
now written as a function of its complex argument. Let \(\mathcal N_S\)
be the positive integers prime to every prime in \(S\). Write
\[
 Z_S^{\rm geom}U(x)=\sum_{n\in\mathcal N_S}n^{-1/2}U(x+\log n),
 \qquad T_lU(x)=U(x+l).
 \tag{MA-MO.44}
\]
The superscript means the actual convergent summation, with no Mellin
continuation or pole removal. The operators \(Z_{a,S},J_{a,S}\) remain
the maximal closed multipliers with symbols
\(A_S(t)=E_S(\sigma+it)\zeta(\sigma+it)\) and \(A_S(t)^{-1}\).
The geometric inverse \(J_S^{\rm geom}\) is the original Euler inverse
on Meyer's first smooth component, obtained on every wider
\(\mathcal H_b\), \(b>1/2\), by its absolutely convergent operator sum.

Use precisely the source orientation of the cutoff:
\(\phi\in C^\infty(\mathbb R)\) is one sufficiently far to the left
and zero sufficiently far to the right. In particular its derivative
is compactly supported and \(\phi\) is nonconstant. Define
\[
                 \ell_\phi(F)=\mathcal M(\phi F)(1/2).
 \tag{MA-MO.45}
\]
Symmetry \(\phi(x)+\phi(-x)=1\) will only be needed for the reflected
second-diagonal observation at the end.

\paragraph{MA-MO-004 (exact geometric and continued comparison).}
For every \(U\in\mathcal E\), the following equality holds as an
equality of smooth functions, and its right side belongs to
\(\mathcal H_a\):
\[
 Z_{a,S}U=Z_S^{\rm geom}U-\kappa_S\mathcal MU(1/2)p_0.
 \tag{MA-MO.46}
\]
If \(F\in\mathcal E_0\) and \(y=Z_S^{\rm geom}F\), then \(y\in
\mathcal E\) is an actual smooth theta-image vector and
\[
                       J_S^{\rm geom}y=J_{a,S}y=F.
 \tag{MA-MO.47}
\]
All the following compositions are defined on these vectors:
\[
 \begin{aligned}
 Q_\phi^{\rm geom}y&=Z_S^{\rm geom}(\phi F),\\
 Q_\phi^{\rm cont}y&=Z_{a,S}(\phi F),\\
 D^{\rm geom}y&=Z_S^{\rm geom}((1-\phi)F)-(1-\phi)y\\
              &=\phi y-Q_\phi^{\rm geom}y,\\
 D^{\rm cont}y&=\phi y-Q_\phi^{\rm cont}y.
 \end{aligned}
 \tag{MA-MO.48}
\]
Their exact comparison is
\[
 \begin{aligned}
 Q_\phi^{\rm cont}y
       &=Q_\phi^{\rm geom}y-\kappa_S\ell_\phi(F)p_0,\\
 D^{\rm cont}y
       &=D^{\rm geom}y+\kappa_S\ell_\phi(F)p_0.
 \end{aligned}
 \tag{MA-MO.49}
\]
Moreover
\[
 Q_\phi^{\rm geom}y\in\mathcal H_a
 \quad\Longleftrightarrow\quad
 D^{\rm geom}y\in\mathcal H_a
 \quad\Longleftrightarrow\quad \ell_\phi(F)=0.
 \tag{MA-MO.50}
\]
When these conditions hold, all four outputs in (MA-MO.48) belong
to \(\mathcal E\), and the continued and geometric outputs agree
exactly. The correction in (MA-MO.49) has rank one in the ambient
smooth-function space. It is not a bounded rank-one operator on
\(\mathcal H_a\), since \(p_0\notin\mathcal H_a\).

\begin{proof}
\emph{The original sum and its polar asymptotic.}
For \(U\in\mathcal E\), the sum (MA-MO.44) and all its derivatives
converge locally uniformly: on a fixed compact interval in \(x\),
the summands are bounded by \(C_N n^{-N-1/2}\) for arbitrary
\(N\). The same estimate with its \(x\)-dependence retained proves
rapid exponential decay, with every derivative, at the right end.
For \(S=\varnothing\), define the even function
\[
              u(t)=\frac{U(\log t)}{2\sqrt t}\quad(t>0),
              \qquad u(0)=0.
 \tag{MA-MO.51}
\]
It is Schwartz and flat at zero. Indeed, each derivative away from
zero is a finite sum of powers of \(t^{-1}\) times derivatives
of \(U(\log t)\), which decay faster than every power at both
ends. Also \(\int_{\mathbb R}u(t)\,dt=\mathcal MU(1/2)\).
The self-dual additive Fourier transform, with kernel
\(e^{-2\pi it\xi}\), and Poisson summation give
\[
 Z_\varnothing^{\rm geom}U(x)
   =\mathcal MU(1/2)e^{-x/2}
       +2e^{-x/2}\sum_{n\ge1}\widehat u(ne^{-x}).
 \tag{MA-MO.52}
\]
The last term and every logarithmic derivative decay faster than
every exponential as \(x\to-\infty\). This follows from the
Schwartz estimates for \(\widehat u\), after choosing their order
arbitrarily large and summing \(n^{-N}\). The condition \(u(0)=0\)
has removed the other Poisson polar term.

Finite deletion is the exact identity
\[
 Z_S^{\rm geom}
    =\prod_{p\in S}(1-p^{-1/2}T_{\log p})
                         Z_\varnothing^{\rm geom}.
 \tag{MA-MO.53}
\]
Absolute local convergence justifies reindexing. Since
\(T_{\log p}p_0=p^{-1/2}p_0\), the coefficient of the left polar
term becomes \(\kappa_S\mathcal MU(1/2)\). Thus the actual sum
belongs to \(\mathcal E\) exactly when this moment vanishes.
Without that cancellation it still lies in the first smooth
component of Meyer: every weight \(b>1/2\) makes its left polar
tail rapidly decreasing, and its right tail is already rapid.

\emph{Identification of the continued multiplier.}
For \(b>1/2\), the original Euler sum is bounded on
\(\mathcal H_b\), by the absolutely convergent norm sum
\(\sum_{n\in\mathcal N_S}n^{-b-1/2}\). Its Mellin transform there is
\[
 E_S(z+1/2)\zeta(z+1/2)\mathcal MU(z),\qquad \Re z=b.
 \tag{MA-MO.54}
\]
Move the inverse Mellin line to \(\Re z=a\). The Mellin transform
of \(U\) decreases faster than every power on this fixed strip,
whereas zeta has polynomial growth away from its sole pole.
Horizontal integrals therefore tend to zero. The only pole crossed
is \(z=1/2\), with residue
\(\kappa_S\mathcal MU(1/2)e^{-x/2}\) after multiplication by
the inverse Mellin kernel. The right-line integral is the
left-line integral plus this residue. This proves (MA-MO.46),
including its sign. Zeros of zeta are not poles in this step,
so no generic-line or zero-free hypothesis is used.

The maximal operator domain also follows directly:
\(V_aU\) is Schwartz and \(A_S\) has polynomial growth, so
\(A_SV_aU\in L^2\). The pole-subtracted function in (MA-MO.46)
decays rapidly at the left; at the right it is a rapidly
decreasing function minus a multiple of \(e^{-x/2}\), which is
square integrable with weight \(e^{2ax}\) because \(a<1/2\).

\emph{The two actual inverse domains.}
If \(F\in\mathcal E_0\), (MA-MO.52)--(MA-MO.53) show that
\(y\in\mathcal E\). Formula (MA-MO.51), now with \(U=F\),
has both value at zero and integral zero, so its theta sum is
the actual smooth rational theta image, as in MA-TR-101 and
MA-MO.39--.40. Finite Euler deletion is a finite linear
combination of translations and preserves that image.
On each \(\mathcal H_b\), \(b>1/2\), the original Euler inverse
is bounded and inverse to the original sum. Consequently it
returns \(F\) on \(y\); this is the source smooth inverse,
not an independently chosen inverse.

On the narrow line, (MA-MO.46) gives
\(V_ay=A_SV_aF\). The line zeros of \(A_S\) are discrete
and have measure zero, while the finite Euler factors never
vanish for \(\sigma>0\). Thus \(A_S^{-1}V_ay=V_aF\) almost
everywhere, and this function is in \(L^2\). This proves
membership in the maximal inverse domain and (MA-MO.47).
Both \(\phi F\) and \((1-\phi)F\) are in \(\mathcal E\).
Their original sums and their continued multipliers are
therefore defined. Linearity of the actual sum proves
(MA-MO.48), and (MA-MO.46) proves (MA-MO.49).

Finally \(p_0\notin\mathcal H_a\), since
\(e^{2ax}|p_0(x)|^2=e^{(2a-1)x}\) is not integrable at
the left end. The continued outputs are in \(\mathcal H_a\)
and \(\phi y\in\mathcal E\). Thus (MA-MO.49) proves the
equivalences (MA-MO.50). If \(\ell_\phi(F)=0\), both
\(\phi F\) and \((1-\phi)F\) have moment zero, so their
actual sums are in \(\mathcal E\) by the proved Poisson
estimate. All remaining assertions follow.
\end{proof}

\paragraph{MA-MO-005 (unboundedness with both polar moments killed).}
For the source cutoff and every finite \(S\), there exist
\(F_M\in\mathcal E\) and actual smooth theta-image vectors
\(y_M=Z_S^{\rm geom}F_M\in\mathcal E\) such that
\[
 \begin{gathered}
 \mathcal MF_M(1/2)=\mathcal M(\phi F_M)(1/2)=0,
       \qquad \sup_M\|y_M\|_{\mathcal H_a}<\infty,\\
 \|Q_\phi^{\rm geom}y_M\|_{\mathcal H_a}\longrightarrow\infty,
 \qquad
 \|D^{\rm geom}y_M\|_{\mathcal H_a}\longrightarrow\infty.
 \end{gathered}
 \tag{MA-MO.55}
\]
Every input, inverse, and final geometric output in this assertion
is a genuine smooth function in \(\mathcal E\). In particular
neither original raw operator has a bounded extension to
\(\mathcal H_a\) agreeing on this pole-free smooth subspace.
The witnesses require neither A nor M, nor any zero on the
chosen line. The published universality input is exactly that
already specified in MA-MO-002.

\begin{proof}
\emph{A bounded moment functional and a fixed correction.}
The orientation of the source cutoff is essential here. Since
\(\phi=0\) sufficiently far right and \(1-2a>0\),
\[
 L_{\phi,a}^2=\int_{\mathbb R}|\phi(x)|^2e^{(1-2a)x}\,dx<\infty,
 \qquad |\ell_\phi(F)|\le L_{\phi,a}\|F\|_{\mathcal H_a}.
 \tag{MA-MO.56}
\]
Choose nonnegative smooth functions \(\eta_L,\eta_R\), each
of integral one, compactly supported respectively in regions
where \(\phi=1\) and \(\phi=0\). Put
\[
 F_*(x)=e^{-x/2}(\eta_L(x)-\eta_R(x)),\qquad
 y_*=Z_S^{\rm geom}F_*.
 \tag{MA-MO.57}
\]
Then \(F_*\in C_c^\infty\), \(\mathcal MF_*(1/2)=0\),
\(\ell_\phi(F_*)=1\), and \(y_*\in\mathcal E\).
The continued output \(Q_\phi^{\rm cont}y_*\) is a fixed
vector of \(\mathcal H_a\). We do not assert that the
separate raw output \(Q_\phi^{\rm geom}y_*\) is in that
Hilbert space; in fact (MA-MO.50) shows it is not.

\emph{Uniformly bounded inverses of the universality inputs.}
Use the rational inputs \(v_M\) of MA-MO.31. On their support,
write \(t=T_M+\xi/M\), \(|\xi|\le1\). By MA-MO.30,
\[
 A_0(t)=e^\xi(1+\eta_M(t)),\qquad
      |\eta_M(t)|<\varepsilon_0<1/100,
      \qquad A_0(t)=\zeta(\sigma+it).
 \tag{MA-MO.58}
\]
The fixed Euler product has bounds
\(0<m_S\le|E_S(\sigma+it)|\le M_S<\infty\).
Normalize the finite-deletion inputs by
\[
 \begin{aligned}
 v_M^S&=\frac{E_S(\sigma+i\,\cdot)\,v_M}
                    {\|E_S(\sigma+i\,\cdot)\,v_M\|_2},
 &u_M&=V_a^{-1}v_M^S,\\
 q_{0,M}&=A_S^{-1}v_M^S
   =\frac{v_M}{A_0\|E_S(\sigma+i\,\cdot)\,v_M\|_2}.
 \end{aligned}
 \tag{MA-MO.59}
\]
They belong to the exact Fourier core of MA-MO-002. Its
norm lower bound, together with (MA-MO.58), gives fixed
constants \(K,c_S,d>0\) such that
\[
 \begin{gathered}
 \|u_M\|_{\mathcal H_a}=1,\qquad
 \|q_{0,M}\|_2\le K:=\frac{e}{m_S(1-\varepsilon_0)},\\
 \|Q_\phi^{\rm cont}u_M\|_{\mathcal H_a}
             \ge R_M:=c_Se^{Md}/M\longrightarrow\infty.
 \end{gathered}
 \tag{MA-MO.60}
\]
The inverse estimate is uniform in the approximating shifts
\(T_M\). It is not an estimate for the inverse on the entire
line or for arbitrary core inputs.

\emph{Graph approximation with inverse control.}
Apply the explicit Gaussian construction MA-MO.36--.41
separately to each \(v_M^S\). It supplies
\(F_{M,\epsilon}\in\mathcal E_0\) and actual smooth images
\(y_{M,\epsilon}=Z_S^{\rm geom}F_{M,\epsilon}\) satisfying
\[
 \begin{aligned}
 V_aF_{M,\epsilon}&\longrightarrow q_{0,M}&&\text{in }L^2,\\
 y_{M,\epsilon}&\longrightarrow u_M&&\text{in }\mathcal H_a,\\
 Q_\phi^{\rm cont}y_{M,\epsilon}
      &\longrightarrow Q_\phi^{\rm cont}u_M&&\text{in }\mathcal H_a.
 \end{aligned}
 \tag{MA-MO.61}
\]
The first convergence follows from the Schwartz convergence
MA-MO.38; it is part of the explicit construction, not a
consequence of a bounded inverse. For each \(M\), choose
\(\epsilon_M>0\) so that all three errors are at most one.
Abbreviate the resulting vectors by \(F'_M,y'_M\). Then
\[
 \|F'_M\|_{\mathcal H_a}\le K+1,\qquad
 \|y'_M\|_{\mathcal H_a}\le2,\qquad
 \|Q_\phi^{\rm cont}y'_M\|_{\mathcal H_a}\ge R_M-1.
 \tag{MA-MO.62}
\]

\emph{An exact bounded correction, before any raw norm comparison.}
Set
\[
 \begin{aligned}
 c_M&=\ell_\phi(F'_M),&|c_M|&\le C_*:=L_{\phi,a}(K+1),\\
 F_M&=F'_M-c_MF_*,&y_M&=y'_M-c_My_*=Z_S^{\rm geom}F_M.
 \end{aligned}
 \tag{MA-MO.63}
\]
Both required moments of \(F_M\) vanish exactly. Moreover
\(\|y_M\|_{\mathcal H_a}\le C_Y:=2+C_*\|y_*\|_{\mathcal H_a}\).
All terms of the following estimate are continued outputs
in the actual Hilbert space, so the triangle inequality is valid:
\[
 \|Q_\phi^{\rm cont}y_M\|_{\mathcal H_a}
 \ge R_M-1-C_*\|Q_\phi^{\rm cont}y_*\|_{\mathcal H_a}
                              \longrightarrow\infty.
 \tag{MA-MO.64}
\]
Now apply MA-MO-004 to the corrected \(F_M\), whose two
moments are zero. Its continued output equals its literal
geometric output, and that output belongs to \(\mathcal E\).
Hence (MA-MO.64) is the desired raw norm divergence. Also
\[
 \|D^{\rm geom}y_M\|_{\mathcal H_a}
    \ge\|Q_\phi^{\rm geom}y_M\|_{\mathcal H_a}
                         -\|\phi\|_\infty C_Y
                                \longrightarrow\infty.
 \tag{MA-MO.65}
\]
The inverse in the original composition is exactly \(F_M\)
by (MA-MO.47). Since
\(\mathcal M((1-\phi)F_M)(1/2)=0\) too, its original
summation output is in \(\mathcal E\). Thus both literal
source expressions in (MA-MO.48) have the asserted domains
and smooth outputs. No polar subtraction remains in the
final witnesses. This proves (MA-MO.55) and rules out a
bounded extension agreeing with either raw operator there.
\end{proof}

\paragraph{Exact original-source domain and the reserved smaller space.}
Meyer (2005), \S5.2, Lemmas~5.3--5.4, and the following
projection definitions, author-manuscript pp.~39--40,
define the full smooth product
\(\mathcal S(\mathcal C)_{>}\oplus\mathcal S(\mathcal C)_{<}\)
and \(p_+^S\) on it. The primary labels are
\texttt{sec:summation}, \texttt{lem:summation}, and
\texttt{lem:Hilm\_p\_embed}. Centering turns the first
component into the space of functions whose products with
\(e^{bx}\) are Schwartz for every \(b>1/2\). Thus the
vectors \((y_M,0)\), with \(y_M\in\mathcal E\), are in
that genuine full smooth product. The actual inverse
\(\mathcal L_S^{-1}\) there returns \(F_M\), and the
first diagonal derived from the defining maps is
\[
       Z_S^{\rm geom}M_{1-\phi}J_S^{\rm geom}-M_{1-\phi}
                         =D^{\rm geom}.
 \tag{MA-MO.66}
\]
It is the first diagonal of \(P_+^S-P_-\). The printed
\texttt{eq:global\_perturbation} has the reversed label on
its left side; the defining-map sign already audited in
MA-MEY-003 is retained here.

A bounded operator on \(\mathcal H_a\oplus\mathcal H_{-a}\)
agreeing with this original matrix on the coordinate core
containing \((y_M,0)\) would have a bounded first diagonal
compression. MA-MO-005 rules out that extension. The source
also uses a different, smaller sum space: by
\texttt{lem:Poisson\_sumpm}, in \S5.3, its centered form is
\[
 \{(y_0,y_1):y_0-y_1=c_0e^{x/2}-c_1e^{-x/2}\}.
 \tag{MA-MO.67}
\]
A general \((y_M,0)\) does not lie in that smaller space.
Consequently the diagonal argument alone does not exclude
cancellation between entries on that subspace with a separately
specified norm, nor does it settle every possible common-space
transport. This distinction is part of the theorem's scope.

If the source cutoff is symmetric, reflection \(Ru(x)=u(-x)\)
is unitary \(\mathcal H_a\to\mathcal H_{-a}\), and the other
diagonal is \(RD^{\rm geom}R\). Reflecting (MA-MO.49) replaces
\(p_0\) by \(e^{x/2}\), and reflecting the two-moment-zero
witnesses proves the corresponding second-diagonal statement.
Again no assertion about the off-diagonal domains is required.

\paragraph{Source and Grothendieck placement.}
The cutoff orientation is stated explicitly in the supplied
Meyer primary text at lines~1403--1407 and~1703--1709.
The fresh source check here covers lines~3790--4088, including
\texttt{eq:Poisson\_summation} and the full-product versus
smaller-sum distinction. Lemma~5.6, pp.~41--42, assumes raw
weight \(\alpha>1\), hence centered \(a>1/2\). Its smoothed
nuclear result is unaffected: MA-MO-005 rules out the stated
raw extension to \(0<a<1/2\), not the source lemma.
The RS smoothed trace-class results and PR actual smoothed
quotient returns retain their proved domains and hypotheses.

The polar correction lies in the same two-dimensional
representation as the source singular term, whose centered
generators are
\[
                            e^{x/2},\quad e^{-x/2}.
 \tag{MA-MO.68}
\]
This is an explicit finite-rank comparison in the ambient
smooth-function category, not a splitting of a Hilbert
extension or of the finite-length spectral jet modules.
The inherited cote~1 nuclear-trace locators, scans~079--080
and~111--112 through ALPOGE\_BRIDGES BR-TR-01/02, require
actual nuclear maps before cyclicity is used. They do not
turn the non-Hilbert polar term into a bounded finite-rank
perturbation. No motivic positivity, new zero-location input,
or impossibility of all alternative transports is asserted.
